% Partie : counterfactuals
% Chapitre : minimal-change-semantics
% Section : true-false

\documentclass[../../../include/open-logic-section]{subfiles}

\begin{document}

\olfileid[fr]{cnt}{min}{tf}

\olsection{Vérité et fausseté des conditionnels contrefactuels}

Un conditionnel contrefactuel $!A \cif !B$ est vrai autrement que par
vacuité lorsque tous les mondes-$!A$ les plus proches sont des
mondes-$!B$, comme le montre la \olref{fig:true}.
\begin{figure}
\begin{center}
\begin{tikzpicture}[scale=.7]\small
  \spheresystem{5}
  \draw (0,0) node {$w$};
  \propositionintersect{0}{4}{40}{3.5}
  {\tikzset{proposition/.append={smooth,tension=1.4}}
  \proposition[shift={(1.6,-1)}]{90}{1}{30}{4}}
  \path (1.5,3) node[below] {$\formula{B}$};
  \path (3,0) node {$\formula{A}$};
\end{tikzpicture}
\caption{Conditionnel contrefactuel vrai autrement que par vacuité}
\ollabel{fig:true}
\end{center}
\end{figure}
Un conditionnel contrefactuel est également vrai en~$w$ si le système
de sphères autour de~$w$ ne contient aucune sphère admettant~$!A$.
Dans ce cas, il est vrai \emph{par vacuité} (voir
la \olref{fig:vacuous}).
\begin{figure}
\begin{center}
\begin{tikzpicture}[scale=.7]\small
  \spheresystem{5}
  \draw (0,0) node {$w$};
  \proposition{0}{6.5}{40}{3.5}
  {\tikzset{proposition/.append={smooth,tension=1.4}}
  \proposition[shift={(1.2,-1)}]{90}{1}{30}{4}}
  \path (1.5,3) node[below] {$\formula{B}$};
  \spherepos{0}{7}{node {$\formula{A}$}}
\end{tikzpicture}
\caption{Conditionnel contrefactuel vrai par vacuité}
\ollabel{fig:vacuous}
\end{center}
\end{figure}

Lorsqu'il existe des mondes-$!A$ les plus proches, un conditionnel
contrefactuel peut être faux de deux façons. Il se peut d'abord que
certains de ces mondes, mais pas tous, soient des mondes-$!B$. Dans ce
cas, $!A \cif \lnot !B$ est également faux (voir
la \olref{fig:false}).
\begin{figure}
\begin{center}
\begin{tikzpicture}[scale=.7]\small
  \spheresystem{5}
  \draw (0,0) node {$w$};
  \propositionintersect{0}{4}{40}{3.5}
  {\tikzset{proposition/.append={smooth,tension=1.4}}
  \proposition[shift={(1.6,0)}]{90}{1}{30}{3}}
  \path (1.5,3) node[below] {$\formula{B}$};
  \path (3,0) node {$\formula{A}$};
\end{tikzpicture}
\caption{Contrefactuel faux, version à conséquent nié fausse}
\ollabel{fig:false}
\end{center}
\end{figure}
Si, au contraire, aucun des mondes-$!A$ les plus proches n'est un
monde-$!B$, alors $!A \cif !B$ est faux. Mais tous les mondes-$!A$ les
plus proches sont des mondes-$\lnot !B$ ; par conséquent,
$!A \cif \lnot !B$ est vrai
(voir la \olref{fig:false-opposite}).
\begin{figure}
\begin{center}
\begin{tikzpicture}[scale=.7]\small
  \spheresystem{5}
  \draw (0,0) node {$w$};
  \propositionintersect{0}{4}{40}{3.5}
  {\tikzset{proposition/.append={smooth,tension=1.4}}
  \proposition[shift={(-1,0)}]{90}{1}{30}{3}}
  \path (-1,3) node[below] {$\formula{B}$};
  \path (3,0) node {$\formula{A}$};
\end{tikzpicture}
\caption{Contrefactuel faux, version à conséquent nié vraie}
\ollabel{fig:false-opposite}
\end{center}
\end{figure}

Contrairement aux conditionnels stricts, les conditionnels
contrefactuels peuvent être contingents. Considérons le modèle à
sphères de la \olref{fig:contingent}. Les mondes-$!A$ les plus proches
de~$u$ sont tous des mondes-$!B$ ; ainsi, $\mSat{M}{!A \cif
  !B}[u]$. Mais certains mondes-$!A$ les plus proches de~$v$ ne sont
pas des mondes-$!B$ ; ainsi, $\mSat/{M}{!A \cif !B}[v]$.
\begin{figure}
\begin{center}
\begin{tikzpicture}[scale=.6]\small
  \spheresystem{7}
  \spheresystem[shift={(0:3.1)},dashed]{4}
  \propositionintersect{45}{4}{40}{4.5}
  \begin{scope}
    \clip \propositionplot{45}{4}{40}{4.5} ;
    \spherelayer[shift={(0:3.1)}]{3}
  \end{scope}
  \proposition[shift={(1.4,-.5)}]{90}{1}{35}{5}
  \path (0,0) node {$u$};
  \path (0:3.1) node {$v$};
  \spherepos{35}{9}{node {$\formula{A}$}}
  \spherepos{80}{9}{node {$\formula{B}$}}
\end{tikzpicture}
\end{center}
\caption{Conditionnel contrefactuel contingent}
\ollabel{fig:contingent}
\end{figure}
\end{document}
